\documentclass[11pt]{article}

\usepackage[margin=1in]{geometry}
\usepackage[T1]{fontenc}
\usepackage{lmodern}
\usepackage{amsmath,amssymb,amsthm}
\usepackage{bm}
\usepackage{booktabs}
\usepackage{graphicx}
\usepackage{siunitx}
\usepackage{tabularx}
\usepackage{array}
\usepackage{algorithm}
\usepackage{algpseudocode}
\usepackage[round,numbers,sort&compress]{natbib}
\usepackage[colorlinks=true,linkcolor=blue,citecolor=blue,urlcolor=blue]{hyperref}

\newcommand{\E}{\mathbb{E}}
\newcommand{\Prob}{\mathbb{P}}
\newcommand{\xx}{\mathbf{x}}
\newcommand{\qq}{\mathbf{q}}
\newcommand{\CVaR}{\operatorname{CVaR}}
\newcommand{\calQ}{\mathcal{Q}}

\title{Appendix A\\
\large The two-stage recourse formulation:\\
from the model to the first numerical solutions}
\author{Sergio Lucero}
\date{\today}

\begin{document}
\maketitle

\section*{Scope}

The thesis solves a \emph{here-and-now} design problem: one rate vector
$\qq\in\calQ$ is chosen before any information about the conductivity field
arrives, and it is that single vector that is certified in
Chapter~\ref{chap:results_2d}. This appendix documents a different model, in
which the operator observes the aquifer part-way through the remediation and may
reset the pumping rates afterwards, and it reports everything that was needed to
compute it: the two-phase discretization, the merged mass ledger, the chained
discrete adjoint, the progressive hedging loop with nonanticipativity on the
first stage only, the verification ladder that fixes each of those pieces, and
the first numerical solutions on a benchmark instance.

The distinction between \emph{model change} and \emph{algorithm change} is the
reason this appendix exists. Progressive hedging is already derived in
Chapter~\ref{chap:optimization}: consensus program, augmented Lagrangian,
multiplier update, adaptive penalty, and the caveats about nonconvexity. What
that chapter cannot do is exhibit a decomposition, because on a here-and-now
program the scenario copies of the decision are all the same nine numbers and the
consensus step is an average. Here the program has $n_w$ shared first-stage rates
and $n_w$ recourse rates \emph{per scenario}, so nonanticipativity is a real
constraint on a strict subset of the variables, and the decomposition has
something to decompose. Section~\ref{sec:appx-model} writes the model,
Sections~\ref{sec:appx-disc}--\ref{sec:appx-ph} the implementation, and
Section~\ref{sec:appx-numbers} the numbers, including a case in which the
two-stage model does \emph{not} pay and the diagnostics that say why.

\paragraph{One convention.} Every number in this appendix comes from
\texttt{configs/opt/benchmark\_160x80.yaml}, a $160\times80$ refinement of the
benchmark (Section~\ref{sec:appx-bench}), not from the $160\times96$ instance
with $\sigma=0.65$, $\ell=15\,\mathrm{m}$ and a $120\,\mathrm{m}$-wide domain
that carries Chapter~\ref{chap:results_2d} and whose certified cost is
$433\,945$. The instance differs in the conductivity statistics and in the
domain, so the costs below are \emph{not} comparable to
Chapter~\ref{chap:results_2d}; they are comparable to each other, which is what a
two-stage comparison needs.

% ---------------------------------------------------------------------------
\section{The model}
\label{sec:appx-model}

The horizon $[0,T]$ is cut at an observation time $t_{\mathrm{obs}}\in(0,T)$.
Rates are held constant within each phase: $\qq^{1}$ on $[0,t_{\mathrm{obs}})$
and $\qq^{2}$ on $[t_{\mathrm{obs}},T]$, with $\qq^{1}\in\mathbb{R}^{n_w}$ chosen
before anything is observed and $\qq^{2}\in\mathbb{R}^{n_w}$ chosen after. The
information available at $t_{\mathrm{obs}}$ is the measured concentration and
head field, in practice a handful of probes; in the discretization below the
second-stage decision is indexed by the scenario, which is the
\emph{anticipative} relaxation discussed in Section~\ref{sec:appx-caveats} and
gives an upper bound on what adaptivity is worth.

With $\xi$ the random log-conductivity, the program is
\begin{equation}
\begin{aligned}
(\mathrm{P}_2)\qquad
\min_{\qq^{1}\in\calQ}\quad
& \E_{\xi}\bigl[J(\qq^{1},\qq^{2}(\xi),\xi)\bigr]
  +\rho_g\Bigl[\CVaR_{\alpha}\bigl(G(\qq^{1},\qq^{2}(\cdot),\cdot)\bigr)\Bigr]_+^{2}
  +\rho_h\Bigl[\CVaR_{\alpha_h}\bigl(-H(\qq^{1},\qq^{2}(\cdot),\cdot)\bigr)\Bigr]_+^{2}
\\
& \qq^{2}(\xi)\in\calQ \quad\text{for every }\xi,
\qquad
G(\qq^{1},\qq^{2}(\xi),\xi)=\max_{t\ge t_0}c(\xx_c,t)-\bar c,
\end{aligned}
\label{eq:appx-p2}
\end{equation}
with the same cost $J$, the same compliance functional $G$, the same drawdown
margin $H$ and the same penalties $\rho_g,\rho_h$ as the here-and-now program
(P) of Chapter~\ref{chap:introduction}. Only two things change: the decision is a
pair $(\qq^{1},\qq^{2}(\xi))$ instead of a single vector, and the cost of pumping
in the second phase is incurred at the second-phase rates.

Two consequences are worth stating before any code.

\paragraph{(P$_2$) contains (P).} Setting $\qq^{2}(\xi)=\qq^{1}$ for every $\xi$
returns a feasible point of $(\mathrm{P}_2)$ with exactly the value of the
here-and-now program. Hence $F_2^\star\le F_1^\star$ and any run that reports
the reverse has a convergence or implementation error, not a modelling insight.
Section~\ref{sec:appx-numbers} uses this monotonicity as a diagnostic and reports
one instance where it is not yet attained.

\paragraph{The value of the extra freedom is measurable.} With $\mathrm{RP}$ the
optimal value of $(\mathrm{P}_2)$, $\mathrm{EEV}$ the value of the
expected-value (mean-field) design re-evaluated on the true ensemble, and
$\mathrm{WS}$ the wait-and-see value of a program solved separately for each
scenario \citep{birge2011},
\begin{equation}
\operatorname{VSS}=\mathrm{EEV}-\mathrm{RP},
\qquad
\operatorname{EVPI}=\mathrm{RP}-\mathrm{WS}.
\label{eq:appx-vss}
\end{equation}
Neither number is available for the here-and-now model, and both are what a
reviewer of a stochastic-programming paper asks for first. The caveat that comes
with them here --- that a mean-field design can be \emph{infeasible} on the
ensemble, in which case $\mathrm{EEV}$ is a lower bound and VSS is not
interpretable --- is enforced in the runner rather than left to the reader
(Section~\ref{sec:appx-numbers}).

% ---------------------------------------------------------------------------
\section{Two phases, one horizon}
\label{sec:appx-disc}

\subsection{Why the simulator had to be split, not rewritten}

The transport solver advances a single rate vector over a fixed horizon: the
advection operator is assembled once from the steady flow field induced by
$\qq$, and the time loop reuses it at every step. A rate change at
$t_{\mathrm{obs}}$ therefore invalidates the operator, not just the initial
condition. The implementation takes the cheapest correct route: run the solver
twice. The first call integrates $[0,t_{\mathrm{obs}})$ at $\qq^{1}$; the second
integrates $[t_{\mathrm{obs}},T]$ at $\qq^{2}$ with the first call's final
concentration as its initial condition, which the solver already accepts (a
\texttt{c\_init} argument, added for the sandbox experiments of
Chapter~\ref{chap:formulation} and never used by the thesis until now).

The split is described by a small tuple: the number of steps in each phase
($n_a=\mathrm{round}(t_{\mathrm{obs}}/\Delta t)$, $n_b=n-n_a$), the phase
durations, and how many source steps fall in each phase. The last entry is a
correctness trap: the source is active for a fixed number of steps at the start
of the horizon, and a naive split restarts it in the second phase, injecting mass
that the single-phase run never injects. The rule implemented is that
unconsumed source steps carry over,
$n^{\mathrm{src}}_b=\max(0,n^{\mathrm{src}}-n_a)$, so a split at step $50$ of a
$100$-step run whose source is active for $10$ steps gives
$n^{\mathrm{src}}_a=10$, $n^{\mathrm{src}}_b=0$. On the benchmark this is the
case that matters: the source runs for $365$ days and the observation is at
$1825$ days.

\subsection{Merging the two results}

Both calls return a full result object; the two-stage cost needs them as one.
Four merge rules, in order of how easy they are to get wrong.

\paragraph{Fields and records.} \texttt{C\_well} and \texttt{C\_record} are
per-step records with one row per step, so they concatenate directly. The
concentration \emph{fields} do not: the second call's first field is the first
call's last field, and concatenating both would duplicate the split state. The
merged array is $\texttt{C\_fields}_a$ followed by $\texttt{C\_fields}_b[1:]$,
and the merged time axis repeats the same treatment.

\paragraph{The mass ledger has three kinds of entries.} This is the rule that
cost a full debugging session, so it is stated flatly: \texttt{total} is the
absolute stored mass at each step, a \emph{state}; \texttt{source},
\texttt{well}, \texttt{outflow} and \texttt{decay} are \emph{cumulative}
integrals from the start of the run. Because the second phase restarts from the
first phase's final state, \texttt{total} concatenates with \emph{no} offset and
the four integrals take the first phase's last value as an offset. Offsetting
\texttt{total} as well double-counts the carried mass and produces a mass-balance
error around $0.79$ on a run that is physically exact to machine precision ---
which is exactly how it was noticed.

\paragraph{Final state and the well ledger.} \texttt{C\_final} is the second
phase's final field, and the scalar \texttt{integrated\_well\_mass} is summed
across the phases and then reset to the merged well integral, so that it stays
consistent with the ledger it summarises.

\paragraph{Verification.} The merged result must reproduce the single-phase
result when $\qq^{2}=\qq^{1}$, exactly, not approximately. It does: the maximum
absolute difference in $C_{\mathrm{final}}$, in $C_{\mathrm{record}}$, in the
well concentration history and in the merged time axis are all $0.0$ (bitwise
identical, since the two calls execute the same floating-point operations), and
the merged mass balance error is $-8.7\times10^{-16}$. Table~\ref{tab:appx-ver}
collects this and the rest of the ladder.

\subsection{The objective, reassembled}

The cost functional is time-additive only in appearance. Pumping energy is a
power times a duration, treatment is a flow-weighted integral of the extracted
concentration, and the residual-mass and clean-up-time terms read the final state.
With two rate blocks the first two terms split:
\begin{align}
J_{\text{pump}} &= c_e\bigl(\textstyle\sum_i q^1_i\,T_a+\sum_i q^2_i\,T_b\bigr),
&
J_{\text{treat}} &= c_w\,\Delta t\Bigl(\sum_{n<n_a} \qq^{1}\!\cdot\! C^{\text{well}}_n
+\sum_{n\ge n_a} \qq^{2}\!\cdot\! C^{\text{well}}_n\Bigr),
\label{eq:appx-cost}
\end{align}
whereas $J_{\text{res}}$ and $J_{\text{time}}$ are evaluated on the merged final
field and the merged violation trace exactly as in the single-phase objective, so
they need no new code but must not be evaluated on either phase alone.

The drawdown margin needs one genuinely new line: the heads are computed from
$\qq^{1}$ and $\qq^{2}$ separately (flow is linear in the rates, so this is two
cheap flow solves), and the margin that binds is the elementwise minimum of the
two head fields at the observation cells. Taking the minimum is what makes the
penalty on $H$ a statement about the whole horizon rather than about its second
half.

% ---------------------------------------------------------------------------
\section{The adjoint across the split}
\label{sec:appx-adj}

The thesis built a discrete adjoint and then did not use it
(Chapter~\ref{chap:results_3d}); the two-stage model is the first place where it
is genuinely needed, because progressive hedging solves $M$ small nonlinear
programs per iteration --- $18$ variables each, but expensive function
evaluations --- and a derivative-free inner solve would multiply the cost of the
whole campaign by an order of magnitude.

Two facts make the chaining exact rather than approximate.

\emph{The step is linear in the concentration.} The transport step has the form
$C_{n+1}=\bigl(A_{\text{adv}}(\mathbf v(\qq))+A_{\text{disp}}+\mathrm{diag}(\sigma_t+\lambda+\textstyle\sum_i q_i w_i)\bigr)^{-1}C_n$
with $\mathbf v$ linear in $\qq$ and $w_i$ the unit-rate well footprints, so the
discrete adjoint of the phase-$B$ objective with respect to its own initial
condition is a single linear solve, and the multiplier that phase $B$ produces at
its first step is exactly the terminal condition phase $A$ needs. Passing it in
is all the coupling required; the temporal factor $\sigma_t$ is applied inside
the recursion in both phases and therefore needs no special handling at the
interface.

\emph{Everything else is phase-local.} The compliance cell index, the well
footprints and the response weights depend on the geometry and the scenario, not
on the phase, so the two sweeps share the same preprocessing and differ only in
the upwind velocities (from $\qq^{1}$ and $\qq^{2}$) and in the terminal cost.
The derivative with respect to $\qq^{2}$ accumulates in phase $B$; the derivative
with respect to $\qq^{1}$ accumulates in phase $A$ \emph{plus} the phase-$B$
contribution, which arrives through the terminal multiplier and needs no separate
term.

\paragraph{The indexing trap.} The merged violation trace is indexed globally, but
the phase-$B$ adjoint array is indexed locally: a violation at global step $n\ge
n_a$ must be injected at row $n-n_a$ of phase $B$'s tile. The first implementation
used the global index and raised an index error only when a violation happened to
land in the second phase --- a bug invisible on the runs where the maximum
violation occurred early. Both the branch and the fix are covered by a finite-
difference test at two different split points.

\paragraph{Verification.} Central differences on a joint direction, with a
relative step of $10^{-4}$, agree with the adjoint to
$2.6\times10^{-9}$ in $J$ and $1.4\times10^{-10}$ in $G$ at a split at step $20$,
and to $2.6\times10^{-9}$ and $1.1\times10^{-11}$ at a split at step $13$
(Table~\ref{tab:appx-ver}). The cost of the gradient is one extra transport solve
per phase and is independent of $n_w$: adding wells costs nothing in the adjoint.

% ---------------------------------------------------------------------------
\section{Progressive hedging on the two-stage model}
\label{sec:appx-ph}

\subsection{What nonanticipativity means here}

The consensus variable is the first-stage vector $\qq^{1}$ only. The
implementation carries a mask that is one on the first $n_w$ coordinates and zero
on the recourse block, and applies it in exactly two places: the multiplier
update $\lambda_m\leftarrow\lambda_m+r(\qq^{m}-\bar\xx)$ prices disagreement on
the shared block alone, and the reported consensus residual is computed on that
block alone. Per-scenario recourse rates are free to differ; that freedom is the
model.

\subsection{The loop}

\begin{algorithm}[t]
\caption{Progressive hedging for $(\mathrm{P}_2)$ with a CVaR chance constraint}
\label{alg:appx-ph}
\begin{algorithmic}[1]
\State $\qq^{m}\gets\text{initial }(\qq^{1}_0,\qq^{1}_0)$ for all $m$;
       $\lambda_m\gets0$; $r\gets r_{\text{auto}}$
\For{$\text{outer}=1,\dots,\text{max\_outer}$}
  \State evaluate the hedged policy on all $M$ scenarios: $J_m,G_m,H_m$
  \State $(\kappa,w)\gets\CVaR_{\alpha}\bigl(G_{1:M}\bigr)$;\quad
         $h\gets2\rho_g\max(\kappa+s,0)$ \Comment{$s\ge0$ is the DRO slack (0 here)}
  \For{$m=1,\dots,M$}
     \State $\qq^{m}\gets\arg\min_{\qq\in\calQ}\;
        J_m(\qq)+h\,w_m G_m(\qq)+\lambda_m^{\!\top}\qq+\tfrac r2\lVert\qq-\bar\xx^{\text{prev}}\rVert^2$
        \Comment{projected gradient, $\le$ \texttt{inner} steps, backtracking}
  \EndFor
  \State $\bar\xx\gets\frac1M\sum_m\qq^{m}$;\quad
         $\lambda_m\gets\lambda_m+r(\qq^{m}-\bar\xx)\odot\mathsf{mask}$;\quad
         project $\bar\xx$ onto $\calQ$
  \State record $\lVert\bar\xx^{\text{prev}}-\bar\xx\rVert$, $\lVert\qq^{m}-\bar\xx\rVert$;
         stop if both are below tolerance
  \State $r\gets\min(r_{\max},r_{\text{growth}}\,r)$
\EndFor
\end{algorithmic}
\end{algorithm}

Three details of the implementation are not in the textbook description of
\citet{rockafellar1991} and are worth recording.

\paragraph{The CVaR term is linearised, not dualised.} The exact augmented
Lagrangian would carry the risk coupling through a shared auxiliary variable or a
dual on the CVaR constraint. Here the weights $w_m$ and the hinge multiplier $h$
are computed once per outer iteration at the consensus point and frozen while the
$M$ subproblems are solved, which makes each subproblem a differentiable
surrogate of the exact one and makes the outer iteration a Gauss--Seidel sweep
over the linearisation. The consequence is stated as a limitation, not hidden:
the fixed point the loop converges to satisfies the coupled optimality
conditions only up to the linearisation error, and the error is largest when the
tail is thin (Section~\ref{sec:appx-caveats}). A Pyomo model that carries the
$M+1$ coupled conditions exactly is included in the implementation
(\texttt{build\_pyomo\_model}) but has not been solved at benchmark scale.

\paragraph{The proximal penalty scales itself.} $r$ is chosen once from a probe
gradient at the initial consensus point,
$r=\mathrm{clip}\bigl(\rho_g\max_i|\partial G/\partial q_i|/\bar q,10^{-6},r_{\max}\bigr)$,
and then grown geometrically. On the benchmark this gives $r\approx8.0\times10^4$
from $\max_i|\partial G/\partial q_i|=4.0$, against $r\approx1.0\times10^5$ for the
here-and-now run. Both numbers matter for reproducing the results: the value of
$r$ decides how fast the copies are pulled together and therefore how many outer
iterations are needed for the recourse block to exploit its freedom.

\paragraph{The subproblem is solved by projected gradient with backtracking.} Each
subproblem minimises over an $18$-dimensional box intersected with a capacity
simplex, with an exact gradient from Section~\ref{sec:appx-adj}. The step starts
at a fixed fraction of the admissible range scaled by the gradient's infinity
norm, is halved on rejection and grown by $1.5$ on acceptance, and stops when the
scaled step falls below $10^{-10}$ of the range or the point is pinned by the
boundary. Six inner steps per subproblem are enough at benchmark scale because
the warm start from the previous iterate is close.

\subsection{Two reporting traps}

\paragraph{The consensus vector is not a policy.} $\bar\xx$ is the average of the
$M$ scenario copies, so its recourse block is the \emph{ensemble mean} of the
second-stage rates. Evaluating $\bar\xx$ as if it were a decision re-collapses
the model to a single vector applied in both phases --- the two-stage run then
looks \emph{worse} than the one-stage run while being, in fact, a different and
easier problem. The diagnostic that caught this was a table where the two-stage
value exceeded the one-stage value on the same scenarios, which
Section~\ref{sec:appx-model} says is impossible. The reported policy is
therefore the hedged vector $(\bar\qq^{1},\qq^{2}_m)$ for scenario $m$, and the
implementation exposes it as one function so that no caller can get it wrong.

\paragraph{Certification needs the recourse re-solved.} The scenario copies
carried by the algorithm are optimal for the \emph{training} scenarios at the
moment the loop stopped, not for a fresh one. Certifying a two-stage policy on
the test ensemble therefore re-solves the recourse block for each test scenario
with the first stage held fixed, which is the deployable operation (observe, then
re-optimise what is left) and is what the numbers in
Section~\ref{sec:appx-numbers} report. The distinction is not cosmetic: on the
benchmark the two-stage policy scores reliability $0.825$ on the training
ensemble with the algorithm's own copies and $0.950$ on the test ensemble with
re-solved recourse.

% ---------------------------------------------------------------------------
\section{Verification}
\label{sec:appx-ver}

Each piece above has a test that fails loudly if it is wrong, on a deliberately
small fixture: a $16\times8$ grid, $40$ steps, two or three scenarios, two wells,
so that a wrong answer is a wrong answer and not a floating-point argument. The
ladder is ordered so that each rung depends only on the ones above it.

\begin{table}[t]
\centering
\small
\begin{tabularx}{\textwidth}{@{}l l X@{}}
\toprule
Rung & What it fixes & How, and the result \\
\midrule
1 & phase split & split at $0$ or $n$ steps must be rejected, and source steps
  must carry over to the second phase: the check raises, and
  $n^{\text{src}}_b=0$ at split $20$, $n^{\text{src}}-3$ at split $3$ \\
2 & merge and parity & run two phases with $\qq^{2}=\qq^{1}$ and compare with the
  single-phase run: $C_{\text{final}}$, $C_{\text{record}}$, the time axis, $J$,
  $G$ and $\tau$ agree to $\Delta=0.0$ \\
3 & mass ledger & the balance identity is re-checked on the merged ledger:
  $\lvert\cdot\rvert=-8.7\times10^{-16}$ \\
4 & chained adjoint & central differences on a joint direction with
  $\delta=10^{-4}$ at two split points: rel.\ $J=2.6\times10^{-9}$, rel.\
  $G=1.4\times10^{-10}$ \\
5 & gradients reach both phases & perturb $\qq^{1}$ and $\qq^{2}$ separately:
  both partial derivatives are nonzero and they differ \\
6 & degenerate limit & with $M=1$, compare against a well-converged direct solve:
  PH is no worse, and $F=\E[J]+\rho_g[\CVaR]_+^{2}$ holds \\
7 & recourse is worth something & fix $\qq^{1}$ and re-optimise the recourse per
  scenario: $\CVaR$ falls $39.6\to2.0$ for $+7.0\%$ in $\E[J]$ \\
\bottomrule
\end{tabularx}
\caption{Verification ladder for the two-stage implementation, on the small
fixture of the test suite. Rungs 1--6 are regression tests in the repository;
rung 7 is a diagnostic reported here.}
\label{tab:appx-ver}
\end{table}

Rung 7 is the one that shows the mechanism working before any optimization is
attempted: with $\qq^{1}=[5,0]$ fixed on the fixture, re-optimising the recourse
block in each scenario improves $\CVaR_{5\%}(G)$ from $39.57$ to $2.02$ while
increasing expected cost by $7.0\%$. That is the trade the two-stage model exists
to make, and the optimization's job is to find where the trade is worth taking.

% ---------------------------------------------------------------------------
\section{First numerical solutions}
\label{sec:appx-numbers}

\subsection{The instance}
\label{sec:appx-bench}

$200\times100$\,m on $160\times80$ cells ($12\,800$ unknowns), log-normal
conductivity with $\sigma=0.5$ and isotropic correlation length
$\ell=20$\,m truncated at $95\%$ of the variance, nine candidate wells on a
hexagonal lattice with $0\le q_i\le50$\,m$^3$/d and
$\sum_i q_i\le200$\,m$^3$/d, a source of $10$\,m$^3$/d at $100$\,mg/L for
$365$\,d, a single compliance point at $(185,50)$ with $\bar c=1.0$\,mg/L and a
late-time window starting at $80\%$ of the horizon, one drawdown observation at
$(100,50)$ with $h_{\min}=-10$\,m, $\Delta t=36.5$\,d and $n=100$ steps
($T=3650$\,d), $M_{\text{train}}=40$, $M_{\text{test}}=80$,
$\alpha=0.05$, $\rho_g=10^6$. The observation time for the two-stage run is
mid-horizon, $t_{\mathrm{obs}}=50$ steps $=1825$\,d.

\subsection{Value of recourse}

\begin{table}[t]
\centering
\small
\begin{tabularx}{\textwidth}{@{}l l l X@{}}
\toprule
 & here-and-now (P) & two-stage (P$_2$) & \\
\midrule
outer iterations & 13 (converged) & 16 (not converged) & tol.\ $2\times10^{-2}$ / $10^{-3}$ \\
scenario evaluations & 7\,257 & 9\,603 & one per scenario per subproblem step \\
wall time & 8\,164\,s (2.27\,h) & 11\,693\,s (3.25\,h) & one core, no parallelism \\
\midrule
train $\E[J]$ & 422\,165 & 391\,099 & $M=40$ \\
train $\CVaR_{5\%}(G)$ & $+0.230$ & $+2.034$ & train tail is pessimistic (see text) \\
train reliability & 0.950 & 0.825 & hedged policy, algorithm's own copies \\
\midrule
test $\E[J]$ & 422\,068 & 393\,432 & $M=80$, unseen seeds \\
test reliability & 77/80 $=0.963$ & 76/80 $=0.950$ & Wilson $[0.895,0.987]$ / $[0.878,0.980]$ \\
test $\CVaR_{5\%}(G)$ & $+0.278$ & $+0.064$ & not resolvable at $M=80$ \\
\bottomrule
\end{tabularx}
\caption{The same instance solved twice. The two-stage policy costs $6.8\%$ less
in expectation on unseen scenarios at one fewer compliant scenario, and requires
$43\%$ more computation. Both satisfy the $0.94$ reliability gate of
Chapter~\ref{chap:results_2d}.}
\label{tab:appx-bench}
\end{table}

Table~\ref{tab:appx-bench} is the result this appendix exists to report. On the
$80$ scenarios neither run has seen, the two-stage policy costs
$393\,432$ against $422\,068$ for the here-and-now policy --- $6.8\%$ of expected
cost, or $28\,600$ in the units of $J$ --- while giving up one compliant scenario
out of eighty. Both pass the gate; the Wilson intervals overlap, so the honest
statement is that the cost difference is resolved and the reliability difference
is not.

\paragraph{What the second stage actually does.} The first-stage rates are
\[
\qq^{1}=[1.62,\ 0.54,\ 0.05,\ 5.72,\ 0.00,\ 0.13,\ 0.00,\ 0.00,\ 0.44],
\qquad \textstyle\sum_i q^1_i=8.5\ \mathrm{m^3/d},
\]
and the scenario-mean recourse rates are
\[
\overline{\qq^{2}}=[1.87,\ 0.60,\ 0.71,\ 0.38,\ 1.89,\ 0.27,\ 0.13,\ 0.01,\ 0.02],
\qquad \textstyle\sum_i \bar q^2_i=5.9\ \mathrm{m^3/d}.
\]
The interpretation is not the one a first reading suggests. The two-stage design
pumps \emph{half} as much in the first phase as the here-and-now design
($8.5$ against $15.6$\,m$^3$/d), and it moves the duty between wells after the
observation: well 4 carries $5.72$\,m$^3$/d in phase $A$ and between $0.32$ and
$1.13$ in phase $B$, while well 5 carries nothing in phase $A$ and between
$1.71$ and $3.93$ in phase $B$. The recourse is a relocation of pumping towards
the wells that the observed plume actually requires, which is the qualitative
behaviour the 2001 formulation was reaching for and could not compute.

\paragraph{The train tail overstates, again.} The training-ensemble CVaR of the
two-stage policy is $+2.03$, an order of magnitude above the here-and-now run's
$+0.23$, and enough to dominate the reported penalised objective
($4.5\times10^6$, of which $4.1\times10^6$ is the penalty). On the test ensemble
the same policy has $\CVaR_{5\%}(G)=+0.064$. This is the phenomenon of
Chapter~\ref{chap:results_2d} reappearing in a new place: a $5\%$ tail of $40$
training scenarios is an average over two realizations, and it is not a property
of the policy. It also means the \emph{reported penalised objective} of a
two-stage run is not comparable across runs the way $\E[J]$ and reliability are,
because it contains a term that is a property of the sample.

\paragraph{Where the two-stage run has not converged.} The one-stage run met its
tolerances at outer iteration $13$ (consensus residual $1.7\times10^{-2}$ against
$2\times10^{-2}$, step $6.6\times10^{-4}$ against $10^{-3}$). The two-stage run
stopped at its budget of $16$ with a consensus residual of
$5.9\times10^{-2}$, three times its tolerance, and was still improving: the
reported value is therefore an upper bound on $F_2^\star$, and the $6.8\%$ gap is
a lower bound on the value of recourse.

\subsection{Where it does not pay, and why}

The small fixture tells the opposite story, and the contrast is the most useful
diagnostic in this appendix. On the fixture ($40$ steps, three scenarios,
$\alpha=0.34$ so $\alpha M\approx1$, every scenario violating by
$G\approx20$--$60$ in violation units), the two-stage loop is not reliably
better than the here-and-now loop and is unstable across budgets:

\begin{center}
\small
\begin{tabularx}{\textwidth}{@{}l l l l X@{}}
\toprule
budget & one-stage $F$ & two-stage $F$ & ratio & two-stage status \\
\midrule
$4\times3$ & 111\,288 & 212\,021 & 1.905 & not converged \\
$12\times8$ & 107\,574 & 110\,212 & 1.025 & not converged \\
$20\times8$ & 107\,569 & 119\,836 & 1.114 & not converged (13 vs 20 outers) \\
\bottomrule
\end{tabularx}
\end{center}

Two mechanisms explain it, and both are properties of the fixture rather than of
the model. With $\alpha M\approx1$ the frozen CVaR weights place essentially the
entire risk burden on a single scenario, so the other subproblems minimise cost
alone and their violations drift; and because every scenario violates deeply,
the exact penalty is large enough that the linearisation error of
Section~\ref{sec:appx-ph} is not small. The unit-test configuration (two
scenarios, tighter capacities, $12\times8$) does show the expected ordering ---
$F_2=108\,015$ against $F_1=115\,301$, a $6.3\%$ improvement, with the
first-stage rates $[0.10,7.60]$ against the one-stage $[10.0,2.60]$ --- which is
why the monotonicity of Section~\ref{sec:appx-model} is a test and not a claim
about magnitude. The benchmark, with $\alpha M=2$ and a tail that is thin but not
degenerate, is where the mechanism pays.

\subsection{Value of information}

On the fixture the runner also computes the standard trio of
\eqref{eq:appx-vss}. The expected-value design is
$\qq_{\text{ev}}=[0,1.5]$ with $\E[J]=100\,677$ and reliability $0$; the
recourse-problem value is $\E[J]=105\,681$; the wait-and-see value is
$104\,708$. So $\mathrm{EVPI}=+973$ and $\mathrm{VSS}=-5004$, i.e.\ negative ---
which is the documented case in which the mean-field design is infeasible on the
ensemble and $\mathrm{EEV}$ is a lower bound on the cost of any feasible policy.
The runner reports the status alongside the number for exactly this reason, and
at benchmark scale the VSS/EVPI block has not been run yet; it is the next
computation, not a result.

% ---------------------------------------------------------------------------
\section{Costs}
\label{sec:appx-cost}

Three numbers govern any budget built on this code.

\paragraph{One scenario evaluation.} On the benchmark instance, a forward plus
adjoint solve at one scenario and one candidate costs $2.0$\,s for the
single-phase objective and $2.2$\,s for the two-phase objective (two transport
calls, chained adjoint) on one core in NumPy. The second phase is cheap because
the flow solves are linear in the rates and are reused.

\paragraph{One outer iteration.} $(1+\texttt{inner})M$ scenario evaluations when
no backtracking occurs, with the observed counts on this instance between $450$
and $750$ --- the excess over $280$ is backtracking and the consensus
re-evaluation. At $2.2$\,s this is $16$--$28$ minutes per outer iteration, which
is why $16$ outer iterations cost $3.25$\,h.

\paragraph{Threads, not processes.} Multi-threaded BLAS is actively harmful here:
the same evaluation takes $5.28$\,s at eight OMP threads and $2.02$\,s at one, a
$2.6\times$ penalty from oversubscription in the vectorised inner loops. All
runners therefore pin the thread count to one and the scaling route is a process
per scenario set, which is not implemented: the sparse factorizations in the
archive are not picklable, so a worker must rebuild its own scenarios and pay one
flow solve per scenario before it starts.

% ---------------------------------------------------------------------------
\section{Limitations}
\label{sec:appx-caveats}

\begin{enumerate}
\item \textbf{The recourse is anticipative.} Indexing $\qq^{2}$ by the scenario
  lets the second stage see the realized field, whereas an operator sees only the
  state at $t_{\mathrm{obs}}$. The measured value of recourse is therefore an
  upper bound, and the implementable object is a decision rule
  $\qq^{2}=\pi(\text{observation})$ --- an affine rule in the observed
  concentrations is the obvious next step, and it is also the setting in which
  the value of the two-stage model would have to be re-measured before it is
  claimed.
\item \textbf{The CVaR coupling is linearised.} Section~\ref{sec:appx-ph}: the
  fixed point is exact for the consensus term and inexact for the risk term, with
  the error largest when $\alpha M\lesssim2$. The exact coupled program is
  written (\texttt{build\_pyomo\_model}: shared first stage, nonanticipativity
  rows, CVaR epigraph with common $\eta$, external function for the PDE) but has
  not been solved with a general-purpose NLP solver, so the two implementations
  have not been compared.
\item \textbf{One instance, one seed, one observation time.} Table~\ref{tab:appx-bench}
  is $M_{\text{test}}=80$ from a single seed base; the Wilson intervals are wide
  and the cost difference is not paired across seeds. The observation time
  ($1825$\,d) is a choice, not an optimization: a sweep over
  $t_{\mathrm{obs}}\in\{0.25,0.5,0.75\}\,T$ is the natural third axis.
\item \textbf{The two-stage run did not converge} within $16$ outer iterations,
  so the $6.8\%$ is a lower bound on the benefit and the cost comparison
  ($2.27$ against $3.25$\,h) is the comparison at unequal convergence.
\item \textbf{Two dimensions only.} The phase-split machinery is written against
  the two-dimensional solver; the three-dimensional solver would need the same
  merge rules over different tensors, and its cost per evaluation
  (Chapter~\ref{chap:results_3d}) is the reason it was not attempted here.
\end{enumerate}

% ---------------------------------------------------------------------------
\section{Reproduction}

\begin{center}
\small
\begin{tabularx}{\textwidth}{@{}l X@{}}
\toprule
File & Role \\
\midrule
\texttt{stoch/recourse.py} & phase split, merged ledger, two-phase objective,
  \texttt{TwoStageOracle}, finite-difference gradient check \\
\texttt{stoch/adjoint.py} & \texttt{two\_phase\_gradients}: chained sweep with a
  terminal multiplier \\
\texttt{stoch/pha.py} & progressive hedging, hedged-policy reporting, wait-and-see,
  VSS/EVPI, recourse re-solve, Pyomo audit model \\
\texttt{SCRIPTS/run\_pha.py} & command-line driver (\texttt{--mode one|two},
  \texttt{--vss}, \texttt{--t-obs-frac}) \\
\texttt{tests/test\_two\_phase.py} & rungs 1--5 of Table~\ref{tab:appx-ver} \\
\texttt{tests/test\_pha.py} & rung 6, VSS non-negativity, degenerate limit \\
\bottomrule
\end{tabularx}
\end{center}

{\small
\begin{verbatim}
# here-and-now (parity) run: 2.27 h on one core
python3 SCRIPTS/run_pha.py --config configs/opt/benchmark_160x80.yaml \
  --mode one --max-outer 20 --inner 6 --out out/research/pha-one

# two-stage run, observation at mid-horizon: 3.25 h on one core
python3 SCRIPTS/run_pha.py --config configs/opt/benchmark_160x80.yaml \
  --mode two --max-outer 16 --inner 6 --t-obs-frac 0.5 \
  --test-recourse 8 --out out/research/pha-two

# verification ladder only (seconds)
python3 tests/run_tests.py
\end{verbatim}
}

Each run writes \texttt{metrics.json}, \texttt{history.jsonl} (one record per
outer iteration: consensus residual, step norm, $\E[J]$, $\CVaR$, reliability,
evaluation count, iterate) and \texttt{run\_config.json} (git SHA, seeds, device,
wall time) into the output directory, so every number in
Table~\ref{tab:appx-bench} can be re-derived from the artifacts rather than from
this text.

\begin{thebibliography}{9}

\bibitem[Birge and Louveaux(2011)]{birge2011}
J.~R. Birge and F.~Louveaux.
\textit{Introduction to Stochastic Programming}.
Springer, New York, 2nd edition, 2011.

\bibitem[Kleywegt et~al.(2002)]{kleywegt2002}
A.~J. Kleywegt, A.~Shapiro, and T.~Homem-de-Mello.
The sample average approximation method for stochastic discrete optimization.
\textit{SIAM Journal on Optimization}, 12(2):479--502, 2002.

\bibitem[Rockafellar and Uryasev(2000)]{rockafellar2000}
R.~T. Rockafellar and S.~Uryasev.
Optimization of conditional value-at-risk.
\textit{Journal of Risk}, 2(3):21--41, 2000.

\bibitem[Rockafellar and Wets(1991)]{rockafellar1991}
R.~T. Rockafellar and R.~J.-B. Wets.
Scenarios and policy aggregation in optimization under uncertainty.
\textit{Mathematics of Operations Research}, 16(1):119--147, 1991.

\bibitem[Shapiro et~al.(2009)]{shapiro2009}
A.~Shapiro, D.~Dentcheva, and A.~Ruszczy\'nski.
\textit{Lectures on Stochastic Programming: Modeling and Theory}.
SIAM, Philadelphia, 2009.

\bibitem[Watson and Woodruff(2011)]{watson2011}
J.-P. Watson and D.~L. Woodruff.
Progressive hedging innovations for a class of stochastic mixed-integer
resource allocation problems.
\textit{Computational Management Science}, 8(4):355--370, 2011.

\end{thebibliography}

\end{document}
